\section{Problem Formulation (MDP)}\label{sec:problem}

\subsection{Environment configuration}\label{sec:env-config}
We use the highway task of \env{HighwayEnv}~\cite{highway-env} in its faster variant \env{highway-fast-v0}. It shares
the road, vehicle dynamics, observation, actions, termination and reward with \env{highway-v0}, but simulates the
physics at 5\,Hz instead of 15\,Hz, checks collisions only for the ego vehicle, and changes the defaults to 3 lanes,
20 vehicles, 30\,s episodes and an initial ego spacing of 1.5 (\env{highway-v0}: 4 lanes, 50 vehicles, 40\,s, 2). At
20 vehicles it ran about 2.5 times faster per step in our measurements, which made five training seeds per
configuration affordable within the 300k-step budget. To check that the coarser physics does not change the task, we
evaluate every model also on \env{highway-v0} with the configuration of \cref{tab:env-config} (including the ego
spacing of 1.5), at 15\,Hz and with collisions between other vehicles enabled (the \emph{precise} test condition,
\cref{sec:eval-protocol}); for the scripted baselines, crash rates differ by at most 3
percentage points and mean speeds by at most 0.12\,m/s between the two variants. \Cref{tab:env-config} lists the
configuration; it matches \code{BASE\_CONFIG} and \code{DENSITIES} in \code{common/config.py}.

\begin{table}[H]
  \centering
  \small
  \caption{Environment configuration.}
  \label{tab:env-config}
  \begin{tabularx}{\linewidth}{@{}lXX@{}}
    \toprule
    Parameter & Value & Note \\
    \midrule
    Environment          & \env{highway-fast-v0} & training and all test conditions except \emph{precise} \\
    Lanes                & 3 & default of \env{highway-fast-v0} \\
    Other vehicles       & 20 & default of \env{highway-fast-v0} \\
    Episode duration     & 30\,s $=$ 30 decision steps & default of \env{highway-fast-v0} \\
    Policy frequency     & 1\,Hz & one decision per second \\
    Simulation frequency & 5\,Hz & default of the fast variant \\
    Traffic density      & low 0.4, medium 0.7, high 1.0 & \code{vehicles\_density}; training at medium \\
    Ego initial speed    & 25\,m/s & set by the environment \\
    \bottomrule
  \end{tabularx}
\end{table}

The densities were chosen on scripted policies before any agent was trained. At medium density the default cruising
policy (idle) crashed in 54\% of the 50 pilot episodes while a policy that only slows down never crashed, so the
trade-off between speed and safety is present; low density is easy and high density is hard but still feasible for the
slow policy. On the 100 test episodes idle crashes in 71\% (\cref{tab:baselines}); the difference lies within
the sampling noise of 50 pilot episodes, whose 95\% interval is about 40--67\%.

\subsubsection{Observation space}
The observation is the \env{Kinematics} observation of \env{HighwayEnv} with default settings: a $5 \times 5$ array with
one row for the ego vehicle and one for each of the four nearest other vehicles within 200\,m ahead (vehicles up to
about two vehicle lengths behind the ego also qualify), sorted by distance. Each row holds the features presence, $x$,
$y$, $v_x$ and $v_y$. The ego row holds the ego's own values (its normalized $x$ saturates at 1); positions and
velocities of the other vehicles are relative to the ego vehicle, and all features are normalized to $[-1, 1]$ and clipped ($x$ by $\pm 200$\,m, $y$ by $\pm 12$\,m,
velocities by $\pm 80$\,m/s). Rows without a vehicle are zero. Both algorithms flatten the array into a 25-dimensional
input without further normalization. With three lanes, the leader in the ego lane is hidden by the four-vehicle limit
in 0\% of the steps at low and medium density and in 2\% at high density.

\subsubsection{Action space}
The action space is \code{DiscreteMetaAction} with five meta-actions, executed by low-level steering and speed
controllers. We verified their semantics in the simulator (\cref{tab:actions}).

\begin{table}[H]
  \centering
  \small
  \caption{Meta-actions and their effect, verified in \env{HighwayEnv}~1.12.1.}
  \label{tab:actions}
  \begin{tabularx}{\linewidth}{@{}clX@{}}
    \toprule
    Index & Action & Effect \\
    \midrule
    0 & \code{LANE\_LEFT}  & change to the adjacent left lane; the lane index changes within one decision step; no effect in the leftmost lane \\
    1 & \code{IDLE}       & keep the lane and the current target speed \\
    2 & \code{LANE\_RIGHT} & change to the adjacent right lane; no effect in the rightmost lane \\
    3 & \code{FASTER}     & raise the target speed to the next of $\{20, 25, 30\}$\,m/s; no effect at 30\,m/s \\
    4 & \code{SLOWER}     & lower the target speed to the previous of $\{20, 25, 30\}$\,m/s; no effect at 20\,m/s \\
    \bottomrule
  \end{tabularx}
\end{table}

Lanes are numbered 0 (leftmost) to 2 (rightmost). The speed controller reaches a new target speed within one to two
decision steps (for example $25.0 \to 29.3 \to 29.9$\,m/s after \code{FASTER}). An action that has no effect is still
accepted, so the agent can choose it, and our metrics count the lane changes that actually happened rather than the
chosen actions.

\subsubsection{Termination}\label{sec:termination}
An episode terminates when the ego vehicle collides with another vehicle, and is truncated when the time limit of
30\,s is reached. Leaving the road is not a termination condition (\code{offroad\_terminal} is off). In training, only
a collision is treated as a terminal state; at a truncation the value of the next state is bootstrapped, because the
observation does not contain the remaining time.

\subsubsection{Default reward}
The default reward of \env{HighwayEnv} combines four components, each in $[0,1]$: a collision indicator $c_t$, a
high-speed term $r^{\mathrm{speed}}_t$ (forward speed mapped linearly from $[20, 30]$\,m/s to $[0,1]$ and clipped), a
right-lane term $r^{\mathrm{right}}_t$ (the index of the target lane divided by the number of lanes minus one, so 1 in
the rightmost lane) and an on-road indicator:
\begin{equation}
  \tilde{r}_t = -c_t + 0.4\, r^{\mathrm{speed}}_t + 0.1\, r^{\mathrm{right}}_t,
  \qquad
  r_t = \frac{\tilde{r}_t + 1}{1.5}\; \mathbb{1}[\text{on road}].
  \label{eq:default-reward}
\end{equation}
The weighted sum is mapped linearly from $[-1, 0.5]$ to $[0, 1]$. The lane-change weight is 0 by default and is not
used by the environment in version 1.12.1. Because of the normalization, a step without collision earns at least
0.667, and even the step of a collision earns a small non-negative reward, so surviving is rewarded in itself;
\cref{sec:reward-configs} explains why the custom rewards do not use this normalization.

\subsection{Formal definition}
We treat the task as a discounted Markov decision process $(\mathcal{S}, \mathcal{A}, P, R, \gam)$ over the
observations, with episodes truncated after $T$ steps. Strictly, the observation is a partial view of the traffic and
the 30\,s limit is not part of it, so the task is partially observable and the Markov property holds only
approximately; this is why the time limit is treated as a truncation rather than a terminal state
(\cref{sec:termination}).
\begin{itemize}
  \item $\mathcal{S}$: the \env{Kinematics} observation, a $5 \times 5$ array with entries in $[-1, 1]$ (25
        features). It is a partial view of the traffic, since only the four nearest vehicles ahead are observed;
  \item $\mathcal{A}$: the five meta-actions of \cref{tab:actions}, $|\mathcal{A}| = 5$;
  \item $P(\state' \mid \state, \action)$: the simulator dynamics and the behavior of the other vehicles (IDM car
        following and MOBIL lane changes with randomized parameters), stochastic through the initial traffic and the
        behavior parameters;
  \item $R(\state, \action, \state')$: the default reward~\eqref{eq:default-reward} or one of the custom rewards of
        \cref{sec:reward}; it depends on the outcome of the step (collision, completed lane change);
  \item $\gam = 0.99$ for both algorithms; $T = 30$ decision steps (60 in the \emph{long} test condition).
\end{itemize}
The agent maximizes the expected discounted return
\begin{equation}
  J(\pol) = \E_{\pol}\!\left[\sum_{t=0}^{T-1} \gam^{t}\, R(\state_t, \action_t, \state_{t+1})\right].
  \label{eq:return}
\end{equation}

\subsection{State of the art}\label{sec:sota-mdp}
\todo{Draft and extend.}
\begin{itemize}
  \item \textbf{Other observation types:} occupancy grids, bird's-eye-view or camera images
        (\cite{highway-env,kiran2022drl}); trade-off between information content and sample efficiency.
  \item \textbf{Action abstractions:} discrete meta-actions (delegating low-level control to controllers)
        vs.\ continuous steering/acceleration control.
\end{itemize}
